% Official Solution (Problem Sheet 26, Exercise 1)
\begin{exercisesolution}[Solution to \cref{exc:polarization_quadratic_forms_series26}]
\label{sol:polarization_quadratic_forms_series26}
Let $K$ be a field with $\operatorname{char}(K) \neq 2$, let $V$ be a vector space over $K$, and let $B \colon V \times V \to K$ be a symmetric bilinear form.
Using the bilinearity of $B$, we expand the quadratic form $q_B(v + w)$:
\begin{align*}
q_B(v + w) &= B(v + w, v + w) \\
&= B(v, v) + B(v, w) + B(w, v) + B(w, w) \\
&= q_B(v) + B(v, w) + B(w, v) + q_B(w).
\end{align*}
Since $B$ is symmetric, $B(w, v) = B(v, w)$, which simplifies the expression to:
\[
q_B(v + w) = q_B(v) + 2 B(v, w) + q_B(w).
\]
Rearranging terms:
\[
2 B(v, w) = q_B(v + w) - q_B(v) - q_B(w).
\]
Since $\operatorname{char}(K) \ne 2$, the scalar $2 \in K$ is invertible. Dividing by $2$ yields the polarization identity:
\[
B(v, w) = \frac{1}{2} \big( q_B(v + w) - q_B(v) - q_B(w) \big). \qedhere
\]
\exinfo{Official Solution (Problem Sheet 26, Exercise 1). Derives the polarization identity expressing symmetric bilinear forms in terms of quadratic forms in characteristic $\ne 2$.}
\end{exercisesolution}

% Official Solution (Problem Sheet 26, Exercise 3)
\begin{exercisesolution}[Solution to \cref{exc:sylvester_criterion_matrix_tests}]
\label{sol:sylvester_criterion_matrix_tests}
We apply Sylvester's Criterion (\cref{thm:sylvesters_criterion}), which states that a real symmetric matrix is positive definite if and only if all its leading principal minors $\Delta_k := \det(M_{1 \le i,j \le k})$ are strictly positive for $1 \le k \le n$.
\begin{enumerate}[label=\textbf{(\alph*)}]
    \item \textbf{Matrix $A$:}
    The first two leading principal minors are:
    \[
    \Delta_1 = 3 > 0, \qquad \Delta_2 = \det \begin{pmatrix} 3 & 3 \\ 3 & 1 \end{pmatrix} = 3(1) - 3(3) = 3 - 9 = -6 < 0.
    \]
    Since $\Delta_2 < 0$, $A$ is \textbf{not positive definite}.

    \item \textbf{Matrix $B$:}
    The leading principal minors are:
    \begin{align*}
    \Delta_1 &= 6 > 0, \\
    \Delta_2 &= \det \begin{pmatrix} 6 & 3 \\ 3 & 7 \end{pmatrix} = 42 - 9 = 33 > 0, \\
    \Delta_3 &= \det(B) = 6(56 - 9) - 3(24 - 12) + 4(9 - 28) = 6(47) - 3(12) + 4(-19) \\
    &= 282 - 36 - 76 = 170 > 0.
    \end{align*}
    Since all leading principal minors are strictly positive ($\Delta_1, \Delta_2, \Delta_3 > 0$), $B$ is \textbf{positive definite}.

    \item \textbf{Matrix $C$:}
    Computing the successive leading principal minors:
    \begin{align*}
    \Delta_1 &= 3 > 0, \\
    \Delta_2 &= \det \begin{pmatrix} 3 & 0 \\ 0 & 6 \end{pmatrix} = 18 > 0, \\
    \Delta_3 &= \det \begin{pmatrix} 3 & 0 & -1 \\ 0 & 6 & 1 \\ -1 & 1 & 8 \end{pmatrix} = 3(48 - 1) - 1(0 - (-6)) = 3(47) - 6 = 141 - 6 = 135 > 0, \\
    \Delta_4 &= \det(C) = 3 \cdot \det\begin{pmatrix} 6 & 1 & 1 \\ 1 & 8 & 2 \\ 1 & 2 & 5 \end{pmatrix} + (-1) \cdot \det\begin{pmatrix} 0 & 6 & 1 \\ -1 & 1 & 2 \\ 0 & 1 & 5 \end{pmatrix} \\
    &= 3\big(6(36) - 1(3) + 1(-6)\big) - \big(30 - 1\big) = 3(207) - 29 = 621 - 29 = 592 > 0.
    \end{align*}
    Since all four leading principal minors are strictly positive, $C$ is \textbf{positive definite}. \qedhere
\end{enumerate}
\exinfo{Official Solution (Problem Sheet 26, Exercise 3). Tests positive definiteness of three matrices of sizes 2, 3, and 4 via leading principal minors.}
\end{exercisesolution}

% Official Solution (Problem Sheet 26, Exercise 4)
\begin{exercisesolution}[Solution to \cref{exc:positive_definite_equivalent_characterizations}]
\label{sol:positive_definite_equivalent_characterizations}
Let $A \in \M_{n \times n}(\mathbb{R})$ be symmetric. We prove the cycle of implications \textbf{(A)} $\implies$ \textbf{(B)} $\implies$ \textbf{(C)} $\implies$ \textbf{(A)}.
\begin{itemize}
    \item \textbf{(A) $\implies$ (B):}
    Let $\lambda$ be an eigenvalue of $A$ with non-zero eigenvector $v \in \mathbb{R}^n \setminus \{0\}$.
    By positive definiteness of $A$:
    \[
    0 < v\transp A v = v\transp (\lambda v) = \lambda (v\transp v) = \lambda \|v\|^2.
    \]
    Since $v \ne 0$, $\|v\|^2 > 0$, and dividing by $\|v\|^2$ gives $\lambda > 0$.
    Thus all eigenvalues of $A$ are strictly positive.

    \item \textbf{(B) $\implies$ (C):}
    Since $A$ is real symmetric, the Spectral Theorem (\cref{thm:spectral_thm_matrices}) ensures there exists an orthogonal matrix $O \in \Orth(n)$ such that:
    \[
    A = O D O\transp, \qquad D = \operatorname{diag}(\lambda_1, \dots, \lambda_n).
    \]
    By hypothesis, each $\lambda_i > 0$. We define $\sqrt{D} := \operatorname{diag}(\sqrt{\lambda_1}, \dots, \sqrt{\lambda_n})$ and set:
    \[
    S := O \sqrt{D} O\transp.
    \]
    Then $S$ is symmetric ($S\transp = (O \sqrt{D} O\transp)\transp = O \sqrt{D}\transp O\transp = O \sqrt{D} O\transp = S$), and:
    \[
    S^2 = (O \sqrt{D} O\transp)(O \sqrt{D} O\transp) = O \sqrt{D} (O\transp O) \sqrt{D} O\transp = O (\sqrt{D})^2 O\transp = O D O\transp = A.
    \]
    Moreover, $\det(S) = \det(O)\det(\sqrt{D})\det(O\transp) = \prod_{i=1}^n \sqrt{\lambda_i} > 0$, so $S$ is invertible.

    \item \textbf{(C) $\implies$ (A):}
    Let $S \in \M_{n \times n}(\mathbb{R})$ be an invertible symmetric matrix such that $S^2 = A$.
    For any non-zero vector $v \in \mathbb{R}^n \setminus \{0\}$:
    \[
    v\transp A v = v\transp S^2 v = v\transp S\transp S v = (Sv)\transp (Sv) = \|Sv\|^2.
    \]
    Since $S$ is invertible and $v \ne 0$, we have $Sv \ne 0$.
    By positive definiteness of the Euclidean norm, $\|Sv\|^2 > 0$, so $v\transp A v > 0$, proving that $A$ is positive definite. \qedhere
\end{itemize}
\exinfo{Official Solution (Problem Sheet 26, Exercise 4). Proves cyclic equivalence between positive definiteness, positive eigenvalues, and positive symmetric square root factorization.}
\end{exercisesolution}
